% Section IV -- Optimization Model
% Draft 1 (22 Nov 2026). Follows docs/fable_prompt.md and docs/fable_prompt_round2.md and the MATLAB code in model/.
\section{Optimization Model}
\label{sec:model}

\subsection{Overview}
Component sizes and the hourly dispatch are optimized together as one linear program (LP), following
the micro-grid LP of the course and the joint sizing--dispatch formulations
of~\cite{chedid2020,leon2023}. The model covers a full year at hourly resolution
($t=1,\dots,T$, $T=8760$, $\Delta t=1$~h), so the winter ski peak, multi-day grid outages, and seasonal
storage needs all appear in one problem. The sizing variables are the PV capacity $P_{pv}$ (kWp), the
battery energy and power ratings $E_b$ (kWh) and $P_b$ (kW), and the diesel capacity $P_d$ (kW). The
hourly variables are the PV power used $p_{pv}(t)$ and curtailed $p_{cur}(t)$, battery charging
$p_{ch}(t)$ and discharging $p_{dis}(t)$, stored energy $e(t)$, grid import $p_g(t)$, diesel output
$P_{dg}(t)$ and committed capacity $P_{\mathrm{com}}(t)$, and unserved critical and non-critical load
$\ell_c(t)$ and $\ell_n(t)$. All variables are non-negative and no integer variables are used.

\subsection{Objective}
The model minimizes the annualized cost of supply $C$,
\begin{equation}
\begin{split}
C = {} & \left(\mathrm{CRF}_{pv}c_{pv} + m_{pv}\right) P_{pv}
 + \left(\mathrm{CRF}_{b} + \mu_b\right) c_b E_b \\
 & + \mathrm{CRF}_{d}\,c_d P_d + \sum_{t=1}^{T} \big[ c_g\, p_g(t) + c_f F(t) \\
 & + m_{d} P_{dg}(t) + c_{deg}\, p_{dis}(t) \big]\Delta t,
\end{split}
\label{eq:obj}
\end{equation}
where $c_{pv}$, $c_b$ and $c_d$ are the capital costs, $m_{pv}$ and $\mu_b$ the fixed O\&M of PV and
battery, $c_g$ the EDL tariff, $c_f$ the diesel price, $F(t)$ the fuel use from~\eqref{eq:fuel},
$m_d$ the diesel O\&M, and $c_{deg}$ a throughput cost on battery discharge that represents
wear~\cite{seger2023}. The existing generators are a sunk cost, so $c_d=0$ in the base case. Each
component is annualized over its own life $n$ with the capital recovery factor
$\mathrm{CRF}(i,n) = i(1+i)^n/[(1+i)^n-1]$ at the real discount rate $i$.

\subsection{Constraints}
\emph{Power balance} in every hour $t$ (time index omitted):
\begin{equation}
p_{pv} + p_{dis} + p_g + P_{dg} + \ell_c + \ell_n = P_L + p_{ch}.
\label{eq:balance}
\end{equation}

\emph{PV.} The usable and curtailed PV power add up to the available output,
$p_{pv}(t) + p_{cur}(t) = a(t)\,P_{pv}$, with $a(t)$ from~\eqref{eq:scale} including snow losses.

\emph{Battery.} The stored energy evolves as
\begin{equation}
e(t+1) = e(t) + \eta_{ch}\,p_{ch}(t)\Delta t - \frac{p_{dis}(t)\Delta t}{\eta_{dis}},
\label{eq:soc}
\end{equation}
with $\eta_{ch}=\eta_{dis}=0.9487$ (90\% round trip) and the cyclic condition $e(T+1)=e(1)$, which prevents the year from
ending on stored energy it did not produce. The limits are
\begin{equation}
\mathrm{SOC}_{\min}E_b \le e(t) \le E_b, \quad p_{ch}(t),\,p_{dis}(t) \le P_b \le 0.5\,E_b.
\label{eq:batlim}
\end{equation}

\emph{Grid.} Import is possible only when the feeder is energized, $p_g(t) \le u_g(t)\,P_{g,\max}$,
where $u_g(t)\in\{0,1\}$ is the measured availability of Section~III and $P_{g,\max}$ the connection
capacity, set to the peak load. No energy is exported in the base case.

\emph{Diesel.} Unit commitment normally needs binary variables. The model uses the committed-capacity
relaxation instead, in which $P_{\mathrm{com}}(t)$ is the capacity of the sets running in hour $t$:
\begin{equation}
0.3\,P_{\mathrm{com}}(t) \le P_{dg}(t) \le P_{\mathrm{com}}(t) \le P_d.
\label{eq:diesel}
\end{equation}
The no-load term $\alpha\,P_{\mathrm{com}}$ of~\eqref{eq:fuel} then charges fuel for every committed kW,
which keeps part-load operation costly while the problem stays linear.

\emph{Reliability.} Critical load must always be served, $\ell_c(t)=0$. Non-critical load may be shed
only up to its hourly value, $\ell_n(t) \le P_L(t) - P_{L,c}(t)$, and within an annual
limit on the loss of power supply probability~\cite{yang2008}:
\begin{equation}
\mathrm{LPSP} = \frac{\sum_t \ell_n(t)}{\sum_t P_L(t)} \le \mathrm{LPSP}_{\max}.
\label{eq:lpsp}
\end{equation}
The base case uses $\mathrm{LPSP}_{\max}=0$. A first run with 1\% showed that the model takes the
whole allowance in the busiest ski week, where shedding costs the club most, so values of 0 to 5\%
are reported only as a sensitivity.

\emph{Lift recovery reserve.} While the chairlift runs, the battery holds the energy needed to bring
passengers back on the recovery drive, $e(t) \ge \mathrm{SOC}_{\min}E_b + E_{\mathrm{rec}}$ at the start and
end of the hour, with $E_{\mathrm{rec}}=30$~kWh.

\subsection{Grid Export (Sensitivity)}
Lebanon's 2023 law on distributed renewable energy allows net metering and peer-to-peer sales, but
the credit terms have not yet been set. Export is therefore modelled only as a sensitivity, with a
variable $p_{ex}(t)$ added to the load side of~\eqref{eq:balance} and limited by
\begin{equation}
p_g(t) + p_{ex}(t) \le u_g(t)\,P_{g,\max}, \qquad p_{ex}(t) \le p_{pv}(t),
\label{eq:export}
\end{equation}
so that only PV is exported and only while the feeder is energized, as anti-islanding protection
requires. Under net metering, exported energy is credited at the EDL tariff up to the annual import,
$\sum_t p_{ex}(t) \le \sum_t p_g(t)$; a second case sells it at a fixed price without that cap.

\subsection{Economic Indicators}
The levelized cost of energy (LCOE) is the annualized cost $C$ divided by the energy served. The project is also assessed over $N=25$ years with real, unlevered cash flows:
\begin{equation}
\mathrm{NPV} = -I_0 + \sum_{y=1}^{N} \frac{S_y - R_y}{(1+i)^y} + \frac{V_N}{(1+i)^N},
\label{eq:npv}
\end{equation}
where $I_0$ is the PV and battery investment, $S_y$ the operating saving against the reference supply
in year $y$ (computed with PV output reduced by 0.7\%/yr), $R_y$ the battery replacement in year~15,
and $V_N$ the residual value of the equipment at the end of its life. The internal rate of return is
the value of $i$ for which $\mathrm{NPV}=0$; because the investor and its financing are not known, the
NPV is also reported as a function of $i$.

\subsection{Scenarios and Implementation}
The \emph{reference} case is the current supply: grid and diesel only, with $P_{pv}=E_b=P_b=0$ and the
same load, grid availability, and reliability limit. Every micro-grid case is compared with the
reference of the same scenario. The base case is solved for tilts of 30, 40, 50 and 60$^\circ$; the
sensitivities change one input at a time (load scenario, irradiance, snow loss, grid availability,
diesel price, battery and PV cost, discount rate, reliability limit, cold-weather battery derating, and
grid export). The robustness of the chosen design is checked by fixing its sizes and re-solving the
dispatch for each weather year from 2014 to 2023.

The model is written in MATLAB and solved with \texttt{linprog} (HiGHS dual simplex). The base case has
96,366 variables and 96,362 constraints and solves in about 6~s. It was verified against a 48-hour case
with a hand-calculated optimum and against an independent Python implementation, which gave identical
costs. The model departs from the course formulation in three ways. It imposes no minimum renewable
share, whereas the course slides require at least 60\%; the share is a result, 72.6\% in the base case
but 58.1\% at a diesel price of 0.80~\$/L. The battery may charge from the grid or the diesel plant as
well as from PV. A full 8,760-hour year is modelled instead of one representative day repeated 365
times.
